\pgfplotsset{tuneaxis/.style={width=7.4cm, height=5.2cm, ymin=0.7, ymax=1.0,
  xmin=0.5, xmax=2.5, xlabel={$c$}, ylabel={\scriptsize śr.\ WR vs kontrola},
  ymajorgrids, tick align=outside}}
\begin{figure}[H]\centering
\begin{subfigure}{0.49\textwidth}\centering
\begin{tikzpicture}\begin{axis}[tuneaxis, title={\small UCT @ klasyk}]
\addplot[only marks,mark=*,mark size=1.1pt,color=blue!60]
  table[x=c,y=value]{figs/tune_uct_classic.dat};
\draw[red,densely dashed] (axis cs:2.23,0.7)--(axis cs:2.23,1.0);
\draw[gray,dotted] (axis cs:1.414,0.7)--(axis cs:1.414,1.0);
\end{axis}\end{tikzpicture}\caption{UCT @ klasyk}\end{subfigure}\hfill
\begin{subfigure}{0.49\textwidth}\centering
\begin{tikzpicture}\begin{axis}[tuneaxis, title={\small UCT @ cooldown}]
\addplot[only marks,mark=*,mark size=1.1pt,color=blue!60]
  table[x=c,y=value]{figs/tune_uct_cooldown.dat};
\draw[red,densely dashed] (axis cs:0.52,0.7)--(axis cs:0.52,1.0);
\draw[gray,dotted] (axis cs:1.414,0.7)--(axis cs:1.414,1.0);
\end{axis}\end{tikzpicture}\caption{UCT @ cooldown}\end{subfigure}
\caption{Zależność celu strojenia od stałej eksploracji $c$ (PDP/slice; każdy punkt to jedna próba
Optuny). Czerwona linia przerywana -- najlepsze $c$, szara kropkowana -- $\sqrt2$ (literatura).
Funkcja jest niemal płaska, więc $c$ jest słabo ograniczone.}
\label{fig:tuning}
\end{figure}
